% Options for packages loaded elsewhere
% Options for packages loaded elsewhere
\PassOptionsToPackage{unicode}{hyperref}
\PassOptionsToPackage{hyphens}{url}
\PassOptionsToPackage{dvipsnames,svgnames,x11names}{xcolor}
%
\documentclass[
  british,
]{article}
\usepackage{xcolor}
\usepackage{amsmath,amssymb}
\setcounter{secnumdepth}{5}
\usepackage{iftex}
\ifPDFTeX
  \usepackage[T1]{fontenc}
  \usepackage[utf8]{inputenc}
  \usepackage{textcomp} % provide euro and other symbols
\else % if luatex or xetex
  \usepackage{unicode-math} % this also loads fontspec
  \defaultfontfeatures{Scale=MatchLowercase}
  \defaultfontfeatures[\rmfamily]{Ligatures=TeX,Scale=1}
\fi
\usepackage{lmodern}
\ifPDFTeX\else
  % xetex/luatex font selection
\fi
% Use upquote if available, for straight quotes in verbatim environments
\IfFileExists{upquote.sty}{\usepackage{upquote}}{}
\IfFileExists{microtype.sty}{% use microtype if available
  \usepackage[]{microtype}
  \UseMicrotypeSet[protrusion]{basicmath} % disable protrusion for tt fonts
}{}
\makeatletter
\@ifundefined{KOMAClassName}{% if non-KOMA class
  \IfFileExists{parskip.sty}{%
    \usepackage{parskip}
  }{% else
    \setlength{\parindent}{0pt}
    \setlength{\parskip}{6pt plus 2pt minus 1pt}}
}{% if KOMA class
  \KOMAoptions{parskip=half}}
\makeatother
% Make \paragraph and \subparagraph free-standing
\makeatletter
\ifx\paragraph\undefined\else
  \let\oldparagraph\paragraph
  \renewcommand{\paragraph}{
    \@ifstar
      \xxxParagraphStar
      \xxxParagraphNoStar
  }
  \newcommand{\xxxParagraphStar}[1]{\oldparagraph*{#1}\mbox{}}
  \newcommand{\xxxParagraphNoStar}[1]{\oldparagraph{#1}\mbox{}}
\fi
\ifx\subparagraph\undefined\else
  \let\oldsubparagraph\subparagraph
  \renewcommand{\subparagraph}{
    \@ifstar
      \xxxSubParagraphStar
      \xxxSubParagraphNoStar
  }
  \newcommand{\xxxSubParagraphStar}[1]{\oldsubparagraph*{#1}\mbox{}}
  \newcommand{\xxxSubParagraphNoStar}[1]{\oldsubparagraph{#1}\mbox{}}
\fi
\makeatother


\usepackage{longtable,booktabs,array}
\newcounter{none} % for unnumbered tables
\usepackage{calc} % for calculating minipage widths
% Correct order of tables after \paragraph or \subparagraph
\usepackage{etoolbox}
\makeatletter
\patchcmd\longtable{\par}{\if@noskipsec\mbox{}\fi\par}{}{}
\makeatother
% Allow footnotes in longtable head/foot
\IfFileExists{footnotehyper.sty}{\usepackage{footnotehyper}}{\usepackage{footnote}}
\makesavenoteenv{longtable}
\usepackage{graphicx}
\makeatletter
\newsavebox\pandoc@box
\newcommand*\pandocbounded[1]{% scales image to fit in text height/width
  \sbox\pandoc@box{#1}%
  \Gscale@div\@tempa{\textheight}{\dimexpr\ht\pandoc@box+\dp\pandoc@box\relax}%
  \Gscale@div\@tempb{\linewidth}{\wd\pandoc@box}%
  \ifdim\@tempb\p@<\@tempa\p@\let\@tempa\@tempb\fi% select the smaller of both
  \ifdim\@tempa\p@<\p@\scalebox{\@tempa}{\usebox\pandoc@box}%
  \else\usebox{\pandoc@box}%
  \fi%
}
% Set default figure placement to htbp
\def\fps@figure{htbp}
\makeatother


% definitions for citeproc citations
\NewDocumentCommand\citeproctext{}{}
\NewDocumentCommand\citeproc{mm}{%
  \begingroup\def\citeproctext{#2}\cite{#1}\endgroup}
\makeatletter
 % allow citations to break across lines
 \let\@cite@ofmt\@firstofone
 % avoid brackets around text for \cite:
 \def\@biblabel#1{}
 \def\@cite#1#2{{#1\if@tempswa , #2\fi}}
\makeatother
\newlength{\cslhangindent}
\setlength{\cslhangindent}{1.5em}
\newlength{\csllabelwidth}
\setlength{\csllabelwidth}{3em}
\newenvironment{CSLReferences}[2] % #1 hanging-indent, #2 entry-spacing
 {\begin{list}{}{%
  \setlength{\itemindent}{0pt}
  \setlength{\leftmargin}{0pt}
  \setlength{\parsep}{0pt}
  % turn on hanging indent if param 1 is 1
  \ifodd #1
   \setlength{\leftmargin}{\cslhangindent}
   \setlength{\itemindent}{-1\cslhangindent}
  \fi
  % set entry spacing
  \setlength{\itemsep}{#2\baselineskip}}}
 {\end{list}}
\usepackage{calc}
\newcommand{\CSLBlock}[1]{\hfill\break\parbox[t]{\linewidth}{\strut\ignorespaces#1\strut}}
\newcommand{\CSLLeftMargin}[1]{\parbox[t]{\csllabelwidth}{\strut#1\strut}}
\newcommand{\CSLRightInline}[1]{\parbox[t]{\linewidth - \csllabelwidth}{\strut#1\strut}}
\newcommand{\CSLIndent}[1]{\hspace{\cslhangindent}#1}

\ifLuaTeX
\usepackage[bidi=basic,shorthands=off]{babel}
\else
\usepackage[bidi=default,shorthands=off]{babel}
\fi
\ifLuaTeX
  \usepackage{selnolig} % disable illegal ligatures
\fi


\setlength{\emergencystretch}{3em} % prevent overfull lines

\providecommand{\tightlist}{%
  \setlength{\itemsep}{0pt}\setlength{\parskip}{0pt}}



 


% Get Math Fonts
\usepackage{fontspec}
\usepackage{unicode-math}
\setmathfont{Latin Modern Math}
\setmathfont[range={\mathscr,\mathbfscr,\mathbb}]{XITSMath-Regular}
[    Extension = .otf,
      BoldFont = XITSMath-Bold
]
\makeatletter
\@ifpackageloaded{caption}{}{\usepackage{caption}}
\AtBeginDocument{%
\ifdefined\contentsname
  \renewcommand*\contentsname{Table of contents}
\else
  \newcommand\contentsname{Table of contents}
\fi
\ifdefined\listfigurename
  \renewcommand*\listfigurename{List of Figures}
\else
  \newcommand\listfigurename{List of Figures}
\fi
\ifdefined\listtablename
  \renewcommand*\listtablename{List of Tables}
\else
  \newcommand\listtablename{List of Tables}
\fi
\ifdefined\figurename
  \renewcommand*\figurename{Figure}
\else
  \newcommand\figurename{Figure}
\fi
\ifdefined\tablename
  \renewcommand*\tablename{Table}
\else
  \newcommand\tablename{Table}
\fi
}
\@ifpackageloaded{float}{}{\usepackage{float}}
\floatstyle{ruled}
\@ifundefined{c@chapter}{\newfloat{codelisting}{h}{lop}}{\newfloat{codelisting}{h}{lop}[chapter]}
\floatname{codelisting}{Listing}
\newcommand*\listoflistings{\listof{codelisting}{List of Listings}}
\usepackage{amsthm}
\theoremstyle{plain}
\newtheorem{lemma}{Lemma}[section]
\theoremstyle{definition}
\newtheorem{definition}{Definition}[section]
\theoremstyle{plain}
\newtheorem{proposition}{Proposition}[section]
\theoremstyle{plain}
\newtheorem{theorem}{Theorem}[section]
\theoremstyle{remark}
\AtBeginDocument{\renewcommand*{\proofname}{Proof}}
\newtheorem*{remark}{Remark}
\newtheorem*{solution}{Solution}
\newtheorem{refremark}{Remark}[section]
\newtheorem{refsolution}{Solution}[section]
\makeatother
\makeatletter
\makeatother
\makeatletter
\@ifpackageloaded{caption}{}{\usepackage{caption}}
\@ifpackageloaded{subcaption}{}{\usepackage{subcaption}}
\makeatother

\newcommand{\AmenableGroup}{{\Gamma}}
\newcommand{\Borel}[1]{{\text{Bor}({#1})}}
\newcommand{\ContinuousFunctionSpace}[1]{{\mathscr{C}({#1})}}
\newcommand{\Dual}[1]{{#1}^{*}}
\newcommand{\Folner}[1][\vphantom{}]{{\Phi}_{#1}}
\newcommand{\Group}{{\Gamma}}
\newcommand{\GroupAction}[2]{T^{#1}{#2}}
\newcommand{\GroupActionPreImage}[2]{(T{#1}){-1}.{#2}}
\newcommand{\GroupElement}{{\gamma}}
\newcommand{\Measure}{{\mu}}
\newcommand{\ProjectionMap}[1][\vphantom{}]{{\pi}_{#1}}
\newcommand{\SigmaAlgebra}[1]{{\mathscr{#1}}}
\newcommand{\EigenAlgebra}{{\SigmaAlgebra{K}}}

\usepackage{bookmark}
\IfFileExists{xurl.sty}{\usepackage{xurl}}{} % add URL line breaks if available
\urlstyle{same}
% fallback for those not using the hyperref driver hyperxmp:
\makeatletter
\@ifundefined{xmpquote}{\newcommand{\xmpquote}[1]{#1}}{}
\makeatother
\hypersetup{
  pdftitle={Progressive Measures},
  pdflang={en-GB},
  colorlinks=true,
  linkcolor={blue},
  filecolor={Maroon},
  citecolor={Blue},
  urlcolor={Blue},
  pdfcreator={LaTeX via pandoc}}


\title{Progressive Measures}
\author{}
\date{2026-09-27}
\begin{document}
\maketitle


\begin{definition}[{Kra et al., 2025, Definition
3.1}]\protect\hypertarget{def-ProgressiveMeasure0}{}\label{def-ProgressiveMeasure0}

Let \((X,T)\) be a topological system and \(k\in\mathbb{N}\). We say
that a probability measure \(\tau\in\mathcal{M}(X^{k+1})\) is
\emph{progressive} if, for all open sets \(U_1,...,U_k\subset X\) with
\[\tau(X\times U_1\times\cdots\times U_k)>0 \] there exist infinitely
many \(n\in\mathbb{N}\) such that
\[\tau((X\times U_1\times\cdots\times U_{k-1}\times U_k)\cap T_{\triangle}^{-n}(U_1\times U_2\times\cdots\times U_k\times X))>0. \]

\end{definition}

\begin{definition}[{cf. Kra et al., 2025, Definition
3.1}]\protect\hypertarget{def-ProgressiveMeasure}{}\label{def-ProgressiveMeasure}

Let \((X,\Group)\) be a topological system and \(k\in\mathbb{N}\). We
say that a probability measure \(\tau\in\mathcal{M}(X^{k+1})\) is
\emph{progressive} if, for all open sets \(U_1,...,U_k\subset X\) with
\[\tau(X\times U_1\times\cdots\times U_k)>0 \] there exist infinitely
many \(\GroupElement\in\Group\) such that
\[\tau((X\times U_1\times\cdots\times U_{k-1}\times U_k)\cap \GroupActionPreImage{\GroupElement}{U_1\times U_2\times\cdots\times U_k\times X})>0. \]

\end{definition}

If \((a,x_1,...,x_k)\) is a \(k+1\)-term Erdős progression where \(U_j\)
is a neighbourhood of \(x_j\) for \(1\leq j\leq k\), then
\[B=(X\times U_1\times\cdots\times U_{k-1}\times U_k)\cap T_{\triangle}^{-n}(U_1\times U_2\times\cdots\times U_k\times X)\]
has \((a,x_1,...,x_k)\in B\) and, if further \(\tau(B)>0\), \(B\)
contains infinitely many other \(k+1\)-term Erdős progressions.

Similarly, if \((a,x_{0...01},...,x_{1...1})\) is a \(k\)-dimensional
Erdős cube where \(U_j\) is a neighbourhood of \(x_j\) for all
\(j\in\{0,1\}^k\), \(c_i:\{0,1\}^k\rightarrow\{0,1\}^k\) such that
\[c_i(j)(n)=\begin{cases}1 & j=n,\\ j(n)&\text{otherwise}, \end{cases} \]
and \$C\_, then
\[B=(X\times U_1\times\cdots\times U_{k-1}\times U_k)\cap T_{\triangle}^{-n}(U_1\times U_2\times\cdots\times U_k\times X)\]

\begin{definition}[{cf. Kra et al., 2025, Definition
3.1}]\protect\hypertarget{def-ProgressiveCubicMeasure}{}\label{def-ProgressiveCubicMeasure}

Let \((X,\Group)\) be a topological system and \(k\in\mathbb{N}\). We
say that a probability measure \(\tau\in\mathcal{M}(X^{k+1})\) is
\emph{progressive} if, for all open sets \(U_1,...,U_k\subset X\) with
\[\tau(X\times U_1\times\cdots\times U_k)>0 \] there exist infinitely
many \(\GroupElement\in\Group\) such that
\[\tau((X\times U_1\times\cdots\times U_{k-1}\times U_k)\cap \GroupAction{\GroupElement}{U_1\times U_2\times\cdots\times U_k\times X})>0. \]

\end{definition}

If \((a,x_1,...,x_k)\) is a \(k+1\)-term Erdős progression where \(U_j\)
is a neighbourhood of \(x_j\) for \(1\leq j\leq k\), then
\[B=(X\times U_1\times\cdots\times U_{k-1}\times U_k)\cap T_{\triangle}^{-n}(U_1\times U_2\times\cdots\times U_k\times X)\]
has \((a,x_1,...,x_k)\in B\) and, if further \(\tau(B)>0\), \(B\)
contains infinitely many other \(k+1\)-term Erdős progressions.

\begin{proposition}[{Kra et al., 2025, Proposition
3.2}]\protect\hypertarget{prp-ProgressiveMeasureErdosExistence}{}\label{prp-ProgressiveMeasureErdosExistence}

Fix \(k\in\mathbb{N}\) and let \((X,T)\) be a topological system. Let
\(a\in X\), \(U_1,...,U_k\subset X\) be open sets, and let
\(\tau\in\mathcal{M}(X^{k+1})\) be a progressive measure satisfying
\(\tau(\{a\}\times X^k)=1\).

If \(\tau(X\times U_1\times\cdots\times U_k)>0\), then there exists an
Erdős progression \((a,x_1,...,x_k)\) with \(x_j\in U_j\) for all
\(1\leq j\leq k\).

\end{proposition}

\begin{proof}
Let \(V=U_1\times\cdots\times U_k\). As \(\tau\) is a progressive
measure, we can find \(c(1)\in\mathbb{N}\) such that
\[\tau((X\times V)\cap T_{\triangle}^{-c(1)}(V\times X))>0. \]
\end{proof}

\section{Constructing Progressive
Measures}\label{constructing-progressive-measures}

\begin{enumerate}
\def\labelenumi{\arabic{enumi}.}
\tightlist
\item
  Check that the system has topological pronilfactors.
\item
  Define a product measure in the \(k-1\)-step pronilfactor that is
  generic with \(\Folner\) and \(a\in\text{gen}(\mu,\Folner)\) and
  structured similarly to the progression.
\item
  Project this measure back to the system by using conditional
  expectations.
\end{enumerate}

\section{Claim seemingly related to Progressive
Measures}\label{claim-seemingly-related-to-progressive-measures}

We will look at the following result to see how progressive measures are
involved and determine whether we can generalise this:

\begin{lemma}[{Kra et al., 2026, Claim
1}]\protect\hypertarget{lem-ProgressiveRecurrence}{}\label{lem-ProgressiveRecurrence}

Let \((X,\Measure,T)\) be an ergodic measure-preserving system where
every measurable eigenfunction is equal \(\Measure\)-almost everywhere
to a topological eigenfunction.Also let \(\lambda\) be a self-joining of
\((X,\Measure,T)\), \(d:X\times X\rightarrow[0,\infty)\) denote a metric
on \(X\), and
\((x,y)\in\text{gen}(\lambda,\Folner)\cap(\text{gen}(\Measure,\Folner)\times\text{gen}(\Measure,\Folner))\cap(X\times\text{supp }\Measure)\)
for a Følner sequence \(\Folner\).

For any clopen set \(E\subset X\) with \(\Measure(E)>0\) and any
measurable set \(Y\subset X\times X\): If
\(\lambda\left((X\times E)\cap Y\right)>0\) then, for every
\(\varepsilon>0\), there exist infinitely many
\(\GroupElement\in\Group\) such that \[
d(T^bx,y)<\varepsilon\qquad\text{and}\qquad\lambda\left((T^{-b}X\times E)\cap Y\right)>0.
\]

\end{lemma}

The proof of this result uses the following results:

\ldots{}

We follow and breakdown the proof of
Lemma~\ref{lem-ProgressiveRecurrence} given by Kra et al.
(\citeproc{ref-kra2026shortproof}{2026}):

\begin{proof}
Let \(f_\text{c}=\mathbb{E}(\boldsymbol{1}_E\mid\EigenAlgebra)\) and
\(f_\text{wm}=\boldsymbol{1}_E-f_\text{c}\). As \begin{align*}
  \mathbb{E}(f_\text{wm}\mid\EigenAlgebra)&=\mathbb{E}(\boldsymbol{1}_E-\mathbb{E}(\boldsymbol{1}_E\mid\EigenAlgebra)\mid\EigenAlgebra)
  \\&=\mathbb{E}(\boldsymbol{1}_E\mid\EigenAlgebra)-\mathbb{E}(\boldsymbol{1}_E\mid\EigenAlgebra)
  \\&=0,
\end{align*} then, by \textbf{?@thm-WeakMixing2}, \(f_\text{wm}\) is
mixing.

By \textbf{?@prp-EigenAlgebraProp2}, \(f_\text{c}>0\) for
\(\Measure\)-almost every \(x\in E\). This gives us \begin{align*}
  0<\lambda\left((X\times E)\cap Y\right)&=\int(\boldsymbol{1}\otimes\boldsymbol{1}_E)\cdot\boldsymbol{1}_Y\ \mathrm{d}\lambda
  \\&=\int(\boldsymbol{1}\otimes f_\text{c})\cdot\boldsymbol{1}_Y\ \mathrm{d}\lambda
  \\&=\eta,
\end{align*} where
\(\eta=\langle(\boldsymbol{1}\otimes f_\text{c}),\boldsymbol{1}_Y \rangle\).

\begin{quote}
Though Kra et al. (\citeproc{ref-kra2026shortproof}{2026}) focuses on
\(b\in\mathbb{N}\) from this point on, we would need to generalise to an
amenable group \(\AmenableGroup\).
\end{quote}

As \(f_\text{c}\in\text{L}^2(X,\EigenAlgebra,\Measure)\) and that
\(\text{L}^2(X,\EigenAlgebra,\Measure)\) is defined to be the closed
subspace of \(\text{L}^2(X,\Borel{X},\Measure)\) spanned by its
measurable eigenfunctions, then \(f_\text{c}\) can be approximated by
eigenfunctions. By the assumption that every measurable eigenfunction in
\((X,\Measure,T)\) is equal \(\Measure\)-almost everywhere to a
topological eigenfunction, there exists \(r\in\mathbb{N}\) and
topological eigenfunctions
\(\chi_1,...,\chi_r\in\ContinuousFunctionSpace(X)\) such that
\(f_\text{c}'=\chi_1+\cdots+\chi_r\) satisfies
\(||f_\text{c}'-f_\text{c}||<\eta/6\). Note, by construction,
\(f_\text{c}'\in\ContinuousFunctionSpace(X)\) too.

Let \(\delta>0\) be sufficiently small such that, for \(1\leq i\leq r\)
and for all \(z\in B_\delta(y)\), we have
\(|\chi_i(y)-\chi_i(z)|<\eta/6r\). As \(\chi_i\) is a topological
eigenfunction, we also have
\end{proof}

\begin{itemize}
\tightlist
\item
  Will have to look at finite dimensional representations for Amenable
  Groups setting instead of eigenfunctions

  \begin{itemize}
  \tightlist
  \item
    For actions of abelian groups, the Kronecker factor is spanned by
    eigenfunctions. For actions of general amenable groups, the
    Kronecker factor is made up of finite-dimensional representations.
  \item
    A finite dimensional representation is a homomorphism from G to U(n)
    the group of nxn unitary matrices.

    \begin{itemize}
    \tightlist
    \item
      When n = 1 this is just a homomorphism to the circle.
    \item
      When also G = Z this is just the choice of an eigenvalue.
    \end{itemize}
  \item
    What we are used to is that the Kronecker factor is made up of
    eigenfunctions.

    \begin{itemize}
    \tightlist
    \item
      More generally, the Kronecker factor is made up of spaces that
      look like K/L where K is a compact group and L is a closed
      subgroup. The action of Gamma on each such space will be
      determined by a homomorphism alpha from G to K. Thus the group
      element g in Gamma acts on the coset kL by alpha(g)kL. This is an
      action by isometries because we can always fix an invariant metric
      on K/L. This is enough to generalize the sentence in the claim
      that you mentioned during our meeting.
    \item
      In this more general world, instead of eigenfunctions chi, we have
      ``matrix coefficients'' which are functions of the form chi(x) =
      \textless{} pi(x) u, v \textgreater{} where u and v are vectors in
      L2(K/L).
    \end{itemize}
  \end{itemize}
\item
  Could use the tensor product representation with its dual
  representation?

  \begin{itemize}
  \tightlist
  \item
    Show \(f_c\otimes\bar{f_c}\in\text{fix}(U\otimes\bar{U})\).

    \begin{itemize}
    \tightlist
    \item
      Or construct \(f_c\) such that this is the case.
    \item
      Or write \(f_c\) in terms of its orthonormal bases as it belongs
      to an irreducible invariant finite-dimensional subspace \(M\). If
      so, \(T^g\) would appear to be a matrix transformation \(A_g\) on
      \(M\) that preserves the inner product. The matrix coefficients
      will be defined by \((T^ge_i\mid e_j)\). \(A_g\) is a unitary
      matrix to conserve measures and \(g\mapsto A_g\) is a homomorphism
      mapping from \(G\) to a compact group of unitary matrices (prove).
    \end{itemize}
  \item
    Define topological function analogue, using \[
    \sum_{i=1}^ne_i\otimes\bar{e_i},
    \] where \(\{e_1,...,e_n\}\) is an orthonormal basis of an invariant
    finite-dimensional subspace \(M\subseteq H\).
  \item
    Use the fact that the linear hull of these topological function
    analogues are dense in \(\text{fix}(U\otimes\bar{U})\) to construct
    a topological function analogue that is ``close'' to \(f_c\).
  \end{itemize}
\end{itemize}

\section{Discrete Spectrum \&
Eigenmatrices}\label{discrete-spectrum-eigenmatrices}

\begin{definition}[{cf. Jamneshan and Kreidler, 2025, Definition
5.1.2}]\protect\hypertarget{def-Discrete-Spectrum}{}\label{def-Discrete-Spectrum}

For a unitary representation
\(U:G\rightarrow\mathscr{U}(\mathcal{L}(X,\Borel{X}))\), we call an
invariant finite-dimensional subspace
\(M\subseteq\mathcal{L}(X,\Borel{X})\) \emph{irreducible} if \(\{0\}\)
and \(M\) are the only invariant linear subspaces contained in \(M\).

\end{definition}

\begin{definition}[{cf. Jamneshan and Kreidler, 2025, Definition 5.1.1
\& Corollary
5.1.5}]\protect\hypertarget{def-Discrete-Spectrum}{}\label{def-Discrete-Spectrum}

Let \(U:G\rightarrow\mathscr{U}(\mathcal{L}(X))\) be a unitary
representation of a group \(G\). Then the closure \[
\mathcal{L}(X,\Borel{X})_\text{ds}:=\overline{\text{lin}}\bigcup\{M\subseteq \mathcal{L}(X,\Borel{X})\mid M\text{ irreducible invariant finite-dimensional subspace}\}\subseteq \mathcal{L}(X)
\] is called the \emph{discrete spectrum part} of \(U\). As \(M_1+M_2\)
is an invariant finite-dimensional subspace for invariant
finite-dimensional subspaces \(M_1,M_2\), then
\(\mathcal{L}(X)\text{ds}\) is always a (closed) linear subspace of
\(\mathcal{L}(X)\).

\end{definition}

As there exists \(\ProjectionMap:X\rightarrow K/L\) where \(K\) is a
compact group, \(L\subseteq K\) is a closed subgroup and \(m\) is an
invariant measure on \(K/L\) such that
\(\text{L}^2(X,\SigmaAlgebra{K},\Measure)\cong\text{L}^2(K/L,m)\) and,
for \(f:K/L\rightarrow\mathbb{C}\),
\(f\circ\ProjectionMap:X\rightarrow\mathbb{C}\).

As \(\text{L}^2(X,\SigmaAlgebra{K},\Measure)\cong\text{L}^2(K/L,m)\), we
know that \(\ProjectionMap^{-1}\) exists and thus there exists
\(\bar{f}_\text{c}\in\text{L}^2(K/L,m)\) such that
\(\bar{f}_\text{c}=f_\text{c}\circ\ProjectionMap^{-1}\).

\begin{theorem}[{Farashahi, 2017, Theorem
4.6}]\protect\hypertarget{thm-Peter-Weyl-decomp}{}\label{thm-Peter-Weyl-decomp}

Let \(H\) be a closed subgroup of a compact group \(G\) and \(\mu\) be
the normalized \(G\)-invariant measure on \(G/H\). The Hilbert space
\(\text{L}^2(G/H,\mu)\) satisfies the following orthogonality
decomposition \[
\text{L}^2(G/H,\mu)=\bigoplus_{[\pi]\in\Dual{G/H}}\epsilon_\pi(G/H),
\] where \(\epsilon_\pi\) denotes the linear span of the matrix elements
of \(\pi\). (\(\pi\) may not be the same as the previous one)

\end{theorem}

Note:

\begin{itemize}
\tightlist
\item
  (\citeproc{ref-Farashahi2017PeterWeyl}{Farashahi, 2017}, Theorem 4.5)
  will be helpful in confirming that every continuous function can be
  approximated equicontinuously with a linear combination of
  ``eigenfunctions''.
\end{itemize}

Let \(\{e_1,...,e_n\}\) be the (continuous) orthonormal basis of a
sufficiently large invariant finite-dimensional subspace
\(M\subseteq\text{L}^2(X,\SigmaAlgebra{K},\Measure)\) such that there
exists \(h\in M\) where \(||f_\text{c}-h||_{\text{L}^2}<\eta/6\) and \[
h=\sum_{i=1}^n \langle h,e_i \rangle e_i.
\]

By the Cauchy-Schwartz Inequality, \[
  \langle h,e_i\rangle\leq ||h||_{\text{L}^2}\cdot||e_i||_{\text{L}^2}=||h||_{\text{L}^2}.
\]

We also have \begin{align*}
U_gh&=\sum_{i=1}^n\langle h,e_i \rangle \sum_{j=1}^n \langle U_ge_i,e_j\rangle e_j
\\&=\sum_{j=1}^n \left(\sum_{i=1}^n\langle h,e_i \rangle \langle U_ge_i,e_j\rangle\right) e_j.
\end{align*} As we have \[
\langle U_ge_i,e_j\rangle = \langle e_i,U_g^{-1}e_j\rangle,
\] then \begin{align*}
U_gh&=\sum_{j=1}^n \left(\sum_{i=1}^n\langle h,e_i \rangle \langle U_ge_i,e_j\rangle\right) e_j
\\&=\sum_{j=1}^n \left\langle \sum_{i=1}^n\langle h,e_i \rangle e_i,U_g^{-1}e_j\right\rangle e_j
\\&=\sum_{j=1}^n \langle h,U_g^{-1}e_j\rangle e_j.
\end{align*}

Let \(\delta>0\) be sufficiently small such that, for \(1\leq i\leq n\),
we have \(|e_i(y)-e_i(z)|<\eta/6n\) whenever \(d(z,y)<\delta\). Applying
\(U_g\) and focusing on each \(1\leq i\leq n\), we want to show that \[
|U_ge_i(y)-U_ge_i(z)|<\eta/6n
\] for all \(g\in G\).

We have \[
U_ge_i=\sum_{j=1}^n \langle U_ge_i,e_j\rangle e_j.
\] For each \(e_j\), we have that

\begin{itemize}
\tightlist
\item
  Use continuity as \(e_j\) is assumed to be continuous
\item
  Orthogonal invariance
\item
  Kronecker factor is uniquely ergodic so every point has dense orbit
\end{itemize}

Below is incorrect at the moment:

, as \(e_1,...,e_n\) is an orthonormal basis and
\(|e_j(y)-e_j(z)|<\eta/6n\) for all \(1\leq j\leq n\), we get
\begin{align*}
|U_ge_i(y)-U_ge_i(z)|^2&= \sum_{j=1}^n |\langle U_ge_i,e_j\rangle (e_j(y)-e_j(z))|^2
\\&=\left(\frac{\eta}{6n}\right)^2\sum_{j=1}^n |\langle U_ge_i,e_j\rangle|^2
\\&=\left(\frac{\eta}{6n}\right)^2.
\end{align*} Thus, \[
|U_ge_i(y)-U_ge_i(z)|<\frac{\eta}{6n}.
\]

\subsection{Rough work}\label{rough-work}

By definition of the dual representation, we have \begin{align*}
\overline{h}(f)&=\langle f,h\rangle
\\&=\langle f,\sum_{i=1}^n \langle h,e_i \rangle e_i\rangle
\\&=\sum_{i=1}^n \overline{\langle h,e_i \rangle} \langle f,e_i\rangle
\\&=\sum_{i=1}^n \overline{\langle h,e_i \rangle} \overline{e_i}(f),
\end{align*} for any \(f\in\mathcal{L}(X)\). Thus, \[
\overline{h}=\sum_{i=1}^n \overline{\langle h,e_i \rangle}\overline{e_i}.
\]

With our tensor product
\(U\otimes\overline{U}:G\rightarrow\mathscr{U}(\mathcal{L}(X)\otimes\mathcal{L}(\mathcal{L}(X)))\)
with the dual representation, we get \begin{align*}
(U_g\otimes\overline{U_g})(h\otimes\overline{h})&=(U_gh)\otimes(\overline{U_gh})
\\&=\sum_{i=1}^n\left( \langle h,e_i \rangle U_ge_i\right)\otimes\left(\overline{\langle h,e_i \rangle}\overline{U_ge_i}\right)
\\&=\sum_{i=1}^n\left(\langle h,e_i \rangle \sum_{j=1}^n \langle U_ge_i,e_j\rangle e_j\right)\otimes\left(\overline{\langle h,e_i \rangle}\sum_{k=1}^n\overline{\langle U_ge_i,e_k\rangle}\overline{e_k}\right)
\end{align*} for all \(g\in G\).

Continuing on, we have \begin{align*}
(U_g\otimes\overline{U_g})(h\otimes\overline{h})&=\sum_{i=1}^n\left(\langle h,e_i \rangle \sum_{j=1}^n \langle U_ge_i,e_j\rangle e_j\right)\otimes\left(\overline{\langle h,e_i \rangle}\sum_{k=1}^n\overline{\langle U_ge_i,e_k\rangle}\overline{e_k}\right)
\\&= \sum_{j=1}^n \sum_{k=1}^n\left(\sum_{i=1}^n\langle h,e_i \rangle \overline{\langle h,e_i \rangle} \langle U_ge_i,e_j\rangle \overline{\langle U_ge_i,e_k\rangle}\right) e_j\otimes\overline{e_k}.
\end{align*}

As \(U_x\) is unitary, then \(U_ge_1,...,U_ge_n\) also forms an
orthonormal basis and we get \[
e_k=\sum_{i=1}^n\langle e_k,U_ge_i\rangle U_ge_i.
\]

We also have \[
\sum_{i=1}^n\langle h,e_i \rangle \overline{\langle h,e_i \rangle} \langle U_ge_i,e_j\rangle \overline{\langle U_ge_i,e_k\rangle}=\sum_{i=1}^n\langle\langle e_i,h \rangle h,e_i \rangle \langle\langle e_k,U_ge_i\rangle U_ge_i,e_j\rangle
\]

As \[
  U_gh=\sum_{i=1}^n \langle h,e_i\rangle U_ge_i=\sum_{j=1}^n \langle h,U_g^{-1}e_j\rangle e_j
\] and \(U_ge_1,...,U_ge_n\) are orthonormal, we get \[
\langle h,e_i\rangle U_ge_i=\sum_{j=1}^n\langle h,U_g^{-1}e_j\rangle \langle e_j,U_ge_i\rangle U_ge_i
\]

Using the Cauchy-Schwartz Inequality, we get \begin{align*}
|\langle h,U_g^{-1}e_i\rangle|&\leq ||h||\cdot||U_g^{-1}e_i||
\\& = ||h||\cdot ||e_i||
\\& = ||h||
\end{align*}

We have \[
U_ge_i=\sum_{j=1}^n \langle U_ge_i,e_j\rangle e_j,
\] and so \[
U_ge_i(y)-U_ge_i(z)=\sum_{j=1}^n \langle U_ge_i,e_j\rangle (e_j(y)-e_j(z)).
\] As \(|e_j(y)-e_j(z)|<\eta/6n\) for all \(1\leq j\leq n\), we get
\begin{align*}
|U_ge_i(y)-U_ge_i(z)|&\leq \sum_{j=1}^n \langle U_ge_i,e_j\rangle |(e_j(y)-e_j(z))|
\\&<\frac{\eta}{6n}\sum_{j=1}^n \langle U_ge_i,e_j\rangle
\end{align*}

\protect\phantomsection\label{refs}
\begin{CSLReferences}{1}{0}
\bibitem[\citeproctext]{ref-Farashahi2017PeterWeyl}
Farashahi, A. G. (2017). `Peter-weyl theorem for homogeneous spaces of
compact groups'. \emph{International Journal of Analysis and
Applications}, 13(1), pp. 22--31. Available at:
\url{https://doaj.org/article/937bd45db3334aed9b0a8a029b1b852e}

\bibitem[\citeproctext]{ref-ISem28}
Jamneshan, A. and Kreidler, H. (2025). \emph{ISem 28: Ergodic structure
theory and applications}.

\bibitem[\citeproctext]{ref-kra2025densityfinitesumstheorem}
Kra, B. et al. (2025). \emph{The density finite sums theorem}. Available
at: \url{https://arxiv.org/abs/2504.06424}

\bibitem[\citeproctext]{ref-kra2026shortproof}
Kra, B. et al. (2026). \emph{A short proof of erdős's \(B+C\)
conjecture}.

\end{CSLReferences}




\end{document}
